\documentclass[10pt,a4paper]{amsart}
\usepackage[margin=19mm]{geometry}
\usepackage{amsmath,amssymb,mathtools}
\usepackage[scr=rsfs,cal=euler]{mathalfa}
\usepackage{xfrac}
\usepackage[unicode,hidelinks,bookmarks=false]{hyperref}
\newcommand{\ie}{i.e.\ }
\newcommand{\cf}{cf.\ }
\newcommand{\resp}{respectively\ }
\newcommand{\Subsection}{\subsection}
\providecommand{\nobreakdash}{\nobreak\hbox{-}}
\setlength{\emergencystretch}{2em}
\multlinegap0pt
\usepackage{bm}
\usepackage{bbm}
\usepackage{dsfont}
\usepackage{xspace}
\usepackage{cancel}
\usepackage{mathtools}
\usepackage{amsthm}
\makeatletter
\@ifundefined{Theorem}{\newtheorem{Theorem}{Theorem}}{}
\@ifundefined{Corollary}{\newtheorem{Corollary}[Theorem]{Corollary}}{}
\@ifundefined{Lemma}{\newtheorem{Lemma}[Theorem]{Lemma}}{}
\makeatother
\newcommand{\CC}{\mathbb{C}}
\newcommand{\FF}{\mathbb{F}}
\title[Zeta functions over curves]{Zeta functions over curves}
\author{F. Pellarin, with an appendix by G. H. Ferraro}
\date{Source: arXiv:2606.08848v1 (CC BY 4.0)}
\begin{document}
\maketitle

\noindent\emph{Source and licence.} Zeta functions over curves. Version arXiv:2606.08848v1 (CC BY 4.0), distributed under the Creative Commons Attribution 4.0 International licence, as recorded in the arXiv licence metadata for this exact version and shown on the abstract page https://arxiv.org/abs/2606.08848v1. Full rights notice: licenses/arxiv-2606.08848-CC-BY-4.0.txt.\par\smallskip

\noindent\textit{Editorial scope.} Counted unit: the Lemma whose source label is \texttt{anderson-motives} in the article (source0.tex lines 1024--1027, source label \texttt{anderson-motives}), delivered together with the article's own complete original proof of it (source0.tex lines 1029--1062, one \texttt{proof} environment of 34 lines). No mathematics is written, repaired or re-derived: statement and proof are the article's own text line for line. Required context retained uncounted: corollary-better-bases (lines 919--921); corollary-better-bases-dual (lines 989--991); iota-iso (lines 1013--1015). Every definition, display and label of those lines is retained unchanged. Omitted from this package and not redistributed: 1--918, 922--988, 992--1012, 1016--1023, 1028--1028, 1063--1833. Nothing inside a retained excerpt is omitted or altered: every retained body is delivered exactly as the article's lines are written, so no omission receipt of any kind accompanies this package. Citations: the retained excerpts carry no citation command at all, so no bibliography entry of the article is transcribed and no reference is redistributed. Reference adaptation: none was needed and none was made; every reference inside the delivered unit resolves inside it, so no unresolved reference or citation is produced. Printed slips kept literal: no printed word, symbol, hypothesis or number of the counted statement or of its proof was altered. Layout and class: the article's own class is \texttt{amsart} with packages including upgreek, amssymb, amsmath, geometry, latexsym, graphics, graphicx, tabularx, shapepar, enumerate, hyperref, stmaryrd, glossaries, color; the wrapper is the lane's pinned \texttt{amsart} editorial wrapper at 10pt on A4, which restates the theorem-like environments the retained text needs and adds the article's own macro definitions that the retained text needs (\texttt{\textbackslash CC}, \texttt{\textbackslash FF}), with extra wrapper packages (bm, bbm, dsfont, xspace, cancel, mathtools). The delivered numbering is the unit's own. Assets: no retained span contains a figure environment, a graphics-inclusion command, a PDF-page inclusion, a table or a TikZ picture; no image file is copied into this package, the package declares no asset dependency and no image file is redistributed with it. Licence: version arXiv:2606.08848v1 of this article is distributed under the Creative Commons Attribution 4.0 International licence, as recorded in the arXiv licence metadata for this exact version and shown on the abstract page https://arxiv.org/abs/2606.08848v1. A full rights notice is delivered at licenses/arxiv-2606.08848-CC-BY-4.0.txt. Characters: every retained excerpt is pure ASCII, so the delivered engine, XeLaTeX with its default Latin Modern text font, reports no missing character.
\begin{Corollary}\label{corollary-better-bases}
Consider an integer $n\geq 0$. The $n+1$ functions $ff^{(1)}\cdots f^{(i-1)}$ for $i=0,\ldots,n$ determine a $\CC_\infty$-basis of $\mathcal{F}_{n+1}$.
\end{Corollary}

\begin{Corollary}\label{corollary-better-bases-dual}
Consider an integer $n\geq 0$. The functions ${f^*}^{(-1)}{f^*}^{(-2)}\cdots {f^*}^{(-i)}$ for $i=0,\ldots,n$ determine a $\CC_\infty$-basis of $\mathcal{G}_{n+1}$. Equivalently the functions ${f^*}{f^*}^{(-1)}\cdots {f^*}^{(1-i)}$ for $i=0,\ldots,n$ determine a $\CC_\infty$-basis of $\mathcal{F}^*_{n+1}$.
\end{Corollary}

\begin{equation}\label{iota-iso}
\CC_\infty[\tau]\xrightarrow{\iota}\mathcal{M}.
\end{equation}

\begin{Lemma}\label{anderson-motives}
the vector space $\mathcal{M}$ carries a structure of right $A$-module commuting with the left $\CC_\infty[\tau]$-module structure.
So it can be seen as a left $(A\otimes_{\FF_q}\CC_\infty)[\tau]$-module where $\tau$ commutes with all the elements $a\otimes1$ with $a\in A$.
\end{Lemma}

\begin{proof}
We explain how to define a right action of $A$ on $\mathcal{M}$. Note that for all $a\in A$, $\deg(a)=d$,
$a\otimes1\in\mathcal{L}(d\infty)\subset\mathcal{L}(V+(d+1-1)\infty)=\mathcal{F}_{d+1}$.
Applying Corollary \ref{corollary-better-bases-dual} we can expand, in unique way,
\begin{equation}\label{decomposition-a}
a\otimes1=\sum_{i=0}^{d}[a]_iff^{(1)}\cdots f^{(i-1)}\in\mathcal{F}_{d+1},
\end{equation}
where the coefficients $[a]_0,\ldots,[a]_d$ are in $\CC_\infty$ with $d=\deg(a)$. Note that $f(\Xi)=0$ while $f^{(i)}$ has no pole in $\Xi$ for $i>0$, so that
$$[a]_0=(a\otimes1)(\Xi)=\Xi(a)=a\in\CC_\infty.$$ Also note that $[a]_d\neq0$. Indeed otherwise, observe the equality of divisors
\begin{equation}\label{interated-shtuka}
ff^{(1)}\cdots f^{(i-1)}=V^{(i)}-V+\sum_{j=0}^{i-1}\Xi^{(j)}-i\infty,\quad i\geq0.
\end{equation}
Assuming by contradiction that $[a]_d=0$ for some $a\in A$ of degree $d$, we would have $a\otimes1\in\mathcal{L}((d-1)\infty)$ against the hypothesis on the degree.
For all $n$,
$$(a\otimes1)\mathcal{F}_{n}\subset\mathcal{F}_{d+n}.$$ the multiplication by $a\otimes1$ is therefore an injective 
$\CC_\infty$-endomorphism of $\mathcal{M}$. Coming back to the isomorphism $\iota$ in \eqref{iota-iso}, multiplying by $a\otimes1$ corresponds to the action of an injective endomorphism $\phi_a$ of $\CC_\infty[\tau]$ that has the property that 
for all $n$, it sends $\mathcal{F}_n$ to $\mathcal{F}_{d+n}$. 
Note that for $b\in A$, $(b\otimes1)^{(i)}=b\otimes1$ for all $i$ (the equality of endomorphisms $\tau(b\otimes1)=(b\otimes1)\tau$ follows from the associativity of $\CC_\infty[\tau]$). Therefore, if $b\otimes1=\sum_{j=0}^e[b]_jf\cdots f^{(j-1)}$ (with $e=\deg(b)$), for all $i$,
$$b\otimes1=(b\otimes1)^{(i)}=\sum_{j=0}^{e}[b]_j^{q^i}f^{(i)}f^{(i+1)}\cdots f^{(i+j-1)}.$$ 
Consider $s_i=ff^{(1)}\cdots f^{(i-1)}$ an element of the basis considered in Corollary \ref{corollary-better-bases} (choosing $s_0=1$). 
We have 
\begin{equation}\label{si}
s_i(b\otimes1)=s_i(b\otimes1)^{(i)}=\sum_{j=0}^{e}[b]_j^{q^i}ff^{(1)}\cdots f^{(i-1)}f^{(i)}f^{(i+1)}\cdots f^{(i+j-1)}
=\sum_{i=0}^{d}[b]_i^{q^j}s_{i+j}.\end{equation}

Define, for all $a\in A$,
\begin{equation}\label{drinfeld-module}
\phi_a:=\sum_{i=0}^{\deg(a)}[a]_i\tau^i\in\CC_\infty[\tau],
\end{equation} for coefficients $[a]_i\in\CC_\infty$. By \eqref{si} and uniqueness,
$\phi_a\phi_b=\phi_{ab}$.
Hence, for $a,b\in A$, $\phi_a$ and $\phi_b$ commute: $$\phi_{a}\phi_b=\phi_{ab}=\phi_{ba}=\phi_b\phi_a.$$
Identifying $\mathcal{M}$ with $\CC_\infty[\tau]$ as above via $\iota$, the left multiplication by the rational function
$a\otimes1$ corresponds to the right multiplication by $\phi_a$. 
\end{proof}

\end{document}
